\abstractPT  % Do NOT modify this line

O crescimento das cidades inteligentes exige cidadãos preparados para compreender e participar nas transformações digitais e ambientais que as caracterizam, em alinhamento com os Objetivos de Desenvolvimento Sustentável 11 e~13 das Nações Unidas. Neste contexto, a missão U!REKA promove a transição para cidades mais sustentáveis e inteligentes, identificando a literacia digital dos cidadãos como um vetor estratégico de mudança. O presente trabalho propõe o \textit{Play4Change}, uma plataforma de aprendizagem adaptativa e gamificada destinada a cidadãos, que integra um agente de inteligência artificial capaz de gerar conteúdo formativo, detetar dificuldades de aprendizagem em tempo real e propor percursos alternativos personalizados. A plataforma compreende um cliente móvel, um portal de administração web e um servidor \textit{stateless}. A aprendizagem é organizada em microcompetências progressivas, reconhecidas através de \textit{badges}. Os percursos adaptativos produzidos são reutilizados para cidadãos com perfis de dificuldade semelhantes, potenciando a eficácia coletiva do sistema.

\begin{description}
  \item[Palavras-chave:] Aprendizagem adaptativa, Cidades inteligentes, Gamificação, Microcompetências, Inteligência artificial, U!REKA, ODS
\end{description}
